% =============================================================================
% NeurIPS 2026 Paper Checklist — FINAL (Task 9 audited pass)
% -----------------------------------------------------------------------------
% Self-contained.  Replace the current inline checklist block in main.tex with
%   % =============================================================================
% NeurIPS 2026 Paper Checklist — FINAL (Task 9 audited pass)
% -----------------------------------------------------------------------------
% Self-contained.  Replace the current inline checklist block in main.tex with
%   % =============================================================================
% NeurIPS 2026 Paper Checklist — FINAL (Task 9 audited pass)
% -----------------------------------------------------------------------------
% Self-contained.  Replace the current inline checklist block in main.tex with
%   % =============================================================================
% NeurIPS 2026 Paper Checklist — FINAL (Task 9 audited pass)
% -----------------------------------------------------------------------------
% Self-contained.  Replace the current inline checklist block in main.tex with
%   \input{sections/checklist}
% (the file lives at paper/sections/checklist.tex so the relative path is
% correct when main.tex is compiled from paper/).
%
% Requirements already present in main.tex:
%   \newcommand{\answerPartial}[1][]{\textcolor{purple}{[Partial] #1}}
%   neurips_2026.sty provides \answerYes, \answerNo, \answerNA.
%
% The matching human-readable audit with page references lives at
%   NEURIPS_CHECKLIST_FINAL.md in the repository root.
% =============================================================================

\section*{NeurIPS Paper Checklist}

\begin{enumerate}

% -- 1. Claims -----------------------------------------------------------------
\item {\bf Claims}
    \begin{enumerate}
    \item Do the main claims made in the abstract and introduction
          accurately reflect the paper's contributions and scope?
    \answerYes{The abstract and Section~\ref{sec:intro} state five
    contributions: (i)~a 73-percentage-point implementation gap across
    7~libraries, (ii)~model-dependent GRPO/PPO preference, (iii)~frontier
    collapse on Nemotron-120B, (iv)~zero-variance fraction (ZVF) as a
    leading diagnostic of GRPO failure, and (v)~reward-trajectory
    proxies for KL-free monitoring.  Empirical support is given in
    Section~\ref{sec:results}: the cross-library gap in
    Section~\ref{sec:main_results}; the algorithm$\times$model
    interaction in Section~\ref{sec:ppo_grpo}; the frontier regime in
    Section~\ref{sec:frontier}; ZVF analysis in
    Section~\ref{sec:zvf-analysis}; and proxy metrics in
    Sections~\ref{sec:length_bias} and~\ref{sec:discuss_kl}.  Evidence is
    drawn from the 44~controlled experiments tabulated in
    Section~\ref{sec:results} and the 32-run Bitter Lesson extension of
    Section~\ref{sec:bitter_lesson}.  The scope is explicitly restricted
    to LoRA fine-tuning on three task families (math, code, tool-use)
    with evaluated scales of 0.6B--235B parameters
    (Section~\ref{sec:setup}); Section~\ref{sec:limitations} records
    that 11 of 14 Tinker runs were interrupted by JWT token expiry and
    4 of 6 Modal runs timed out, so reported numbers reflect the
    \emph{available} data and partial runs are marked $\dagger$.}
    \end{enumerate}

% -- 2. Limitations ------------------------------------------------------------
\item {\bf Limitations}
    \begin{enumerate}
    \item Does the paper discuss the limitations of the work performed by
          the authors?
    \answerYes{Section~\ref{sec:limitations} enumerates the limitations
    in detail:
    \begin{itemize}
      \item \textbf{Platform opacity.} Tinker API is a closed, serverless
            platform; GPU type, driver version, CUDA runtime, and
            scheduler are undisclosed, so Tinker-derived results cannot be
            independently reproduced without Tinker access
            (Section~\ref{sec:infra_failures}).
      \item \textbf{Infrastructure failures.} 11~of 14 Tinker runs were
            interrupted by JWT token expiry; 4~of 6 Modal runs hit
            timeouts or gradient-norm blow-ups.  Specific failed runs are
            named in Section~\ref{sec:infra_failures}.
      \item \textbf{Single-seed Tinker runs.} Cost constraints precluded
            multi-seed replication on Tinker; these results are treated
            as descriptive only (Section~\ref{sec:method_limits}).
      \item \textbf{LoRA-only evaluation.} Full fine-tuning and
            quantisation-aware training are not evaluated
            (Section~\ref{sec:method_limits}).
      \item \textbf{Train-set reward metric.} Primary Tinker metrics are
            training rewards, not held-out accuracy, with held-out
            evaluation only completed for the Modal Qwen3-32B / Qwen3.5-27B
            arms (Section~\ref{sec:method_limits}).
      \item \textbf{Proxy metrics in place of direct KL.}
            Section~\ref{sec:discuss_kl} discusses the limits of
            trajectory-based stability indices as substitutes for direct
            KL divergence.
    \end{itemize}}
    \end{enumerate}

% -- 3. Theory ------------------------------------------------------------------
\item {\bf Theory Assumptions and Proofs}
    \begin{enumerate}
    \item For each theoretical result, does the paper provide the full
          set of assumptions and a complete (or correct) proof?
    \answerNA{This is an empirical benchmarking paper; no formal theorems
    are claimed.  GRPO and PPO objectives stated in
    Section~\ref{sec:benchmark} are background definitions, not original
    theoretical results.}
    \end{enumerate}

% -- 4. Reproducibility ---------------------------------------------------------
\item {\bf Experimental Result Reproducibility}
    \begin{enumerate}
    \item Does the paper fully disclose all the information needed to
          reproduce the main experimental results of the paper to the
          extent that it affects the main claims and/or conclusions of
          the paper?
    \answerPartial{Reproducibility is \emph{partial by design} due to
    platform heterogeneity.
    \begin{itemize}
      \item \textbf{Modal experiments (fully reproducible).} Complete
            source code, a pinned \texttt{Dockerfile}, centralised seed
            management (\texttt{utils/seed.py}), and step-by-step
            commands are documented in \texttt{REPRODUCE.md} at
            \url{https://github.com/arvindcr4/tinker-rl-lab}.  Modal
            jobs run on explicitly provisioned NVIDIA H100~SXM5 (80\,GB)
            workers; TRL baselines run on NVIDIA L4 (24\,GB).  All
            Modal/TRL arms use 5~seeds ($s \in \{42, 123, 456, 789, 1024\}$)
            with mean $\pm$ SE and 95\,\% bootstrap CIs via
            \texttt{rliable}.
      \item \textbf{Tinker API experiments (not fully reproducible).}
            Tinker is closed-source with serverless GPU dispatch; GPU
            type, driver version, and scheduling policy are not exposed.
            We release our exact API scripts and configuration files, but
            independent replication requires a Tinker account and is
            subject to the same JWT expiry issues we observed
            (Section~\ref{sec:infra_failures}).  All Tinker-derived
            numbers in figures and tables are flagged with a ``$\dagger$''.
    \end{itemize}
    We claim \emph{Partial} rather than \emph{Yes}: the Tinker subset is
    irreducibly platform-dependent.}
    \end{enumerate}

% -- 5. Open access -------------------------------------------------------------
\item {\bf Open access to data and code}
    \begin{enumerate}
    \item Does the paper provide open access to the data and code, with
          sufficient instructions to faithfully reproduce the main
          experimental results, as described above?
    \answerYes{All code is publicly released under the
    \textbf{Apache License~2.0} at
    \url{https://github.com/arvindcr4/tinker-rl-lab} (mirrored at
    \url{https://github.com/pes-llm-research/tinker-rl-lab}).  LoRA
    adapter checkpoints from the successful Modal experiments are on
    HuggingFace Hub under
    \url{https://huggingface.co/arvindcr4/tinker-rl-bench-*} with model
    cards generated from \texttt{huggingface/MODEL\_CARD\_TEMPLATE.md}.
    All experiment runs are logged publicly at
    \url{https://wandb.ai/arvindcr4-pes-university/tinker-rl-lab-world-class};
    per-step reward, loss, gradient-norm, and (where available) KL
    metrics are visible without login.  Training datasets (GSM8K,
    HumanEval, NoRobots, synthetic tool-use) are public and cited in
    Table~\ref{tab:tasks}.  Tinker API credentials are \emph{not}
    included for security reasons; \texttt{README.md} explains how to
    obtain a Tinker API key.}
    \end{enumerate}

% -- 6. Experimental setting ---------------------------------------------------
\item {\bf Experimental Setting/Details}
    \begin{enumerate}
    \item Does the paper specify all the training and test details
          (e.g., data splits, hyperparameters, how they were chosen,
          type of optimizer) necessary to understand the results?
    \answerYes{Section~\ref{sec:setup} gives models, data splits, LoRA
    configuration (rank~32), optimiser (Adam, $\beta_1{=}0.9$,
    $\beta_2{=}0.95$, $\epsilon{=}10^{-8}$, learning rate $10^{-4}$),
    and the 5-seed Modal protocol.  Table~\ref{tab:hyperparams} maps
    these hyperparameters across all 7~libraries.
    Appendix~\ref{app:hyperparams} reports sweep ranges.  Per-task
    YAML configurations are version-controlled under
    \texttt{atropos/configs/}.  GSM8K uses the standard
    train/test partitions; HumanEval uses the full pass@1 suite.
    Tinker hyperparameters are released as API call scripts; the only
    unavailable information is the Tinker platform's internal hardware.}
    \end{enumerate}

% -- 7. Statistical significance -----------------------------------------------
\item {\bf Experiment Statistical Significance}
    \begin{enumerate}
    \item Does the paper report error bars suitably and correctly
          defined or other appropriate information about the statistical
          significance of the experiments?
    \answerPartial{For Modal/TRL arms with 5~seeds we report mean
    $\pm$ SE with 95\,\% bootstrap confidence intervals via
    \texttt{rliable}~\cite{agarwal2021deep}; a full power analysis
    (Table~\ref{tab:power-analysis}) and Benjamini--Hochberg-adjusted
    pairwise significance (Table~\ref{tab:bh-pvalues}) are reported in
    Section~\ref{sec:setup}.  The TRL-GRPO reference characterisation
    (Qwen2.5-0.5B, 5~seeds) reports $\bar{x}=0.734$,
    $\sigma=0.0703$, IQM $=0.747$, bootstrap 95\,\% CI $[0.672, 0.782]$.
    Tinker API experiments are single-seed (cost-constrained) and are
    explicitly labelled so in Sections~\ref{sec:main_results}
    and~\ref{sec:method_limits}; no statistical significance is claimed
    for Tinker-only comparisons.  We follow the recommendations of
    \citet{colas2019hitchhiker} and \citet{jordan2024benchmarking} on
    the limits of single-seed evaluation.  Because a meaningful subset
    of headline results is Tinker-only, the honest answer is
    \emph{Partial}.}
    \end{enumerate}

% -- 8. Compute resources ------------------------------------------------------
\item {\bf Experiments Compute Resources}
    \begin{enumerate}
    \item For each experiment, does the paper provide sufficient
          information on the computer resources (type of compute
          workers, memory, time of execution) needed to reproduce the
          experiments?
    \answerPartial{Section~4.4 and Appendix~\ref{app:compute} list GPU
    types and estimated GPU-hours for each experiment group; the
    complete per-experiment breakdown with wall-clock time and
    estimated cost is in \texttt{COMPUTE.md}.
    \begin{itemize}
      \item \textbf{Modal experiments.} Fully specified:
            H100~SXM5 (80\,GB) for recent runs and L4 (24\,GB) for the
            TRL reference sweep.  Wall-clock and GPU-hour estimates are
            reported.
      \item \textbf{Tinker API experiments.} Serverless dispatch masks
            the exact GPU type; GPU-hour figures are inferred from
            billing records and are flagged as approximate.
    \end{itemize}
    Total project cost is approximately 1{,}200 A100-equivalent GPU-hours
    across both platforms.  \emph{Partial} reflects the Tinker-side
    hardware opacity, which is a platform property we cannot resolve.}
    \end{enumerate}

% -- 9. Code of Ethics ---------------------------------------------------------
\item {\bf Code Of Ethics}
    \begin{enumerate}
    \item Does the research conducted in the paper conform, in every
          respect, with the NeurIPS Code of Ethics?
    \answerYes{The authors have reviewed the NeurIPS Code of Ethics and
    confirm compliance.  This work involves no human subjects, no
    private or sensitive data, and no models trained on proprietary
    corpora.  All training datasets are public research datasets
    (GSM8K, HumanEval, NoRobots, synthetic tool-use).  Broader societal
    impacts---including reward hacking, alignment failure modes of
    RLHF, and equity of compute access---are discussed in
    Section~\ref{sec:impact} and in \texttt{ethics\_statement.tex}.}
    \end{enumerate}

% -- 10. Broader Impacts -------------------------------------------------------
\item {\bf Broader Impacts}
    \begin{enumerate}
    \item Does the paper discuss both potential positive societal
          impacts and negative societal impacts of the work performed?
    \answerYes{Section~\ref{sec:impact} (plus
    \texttt{ethics\_statement.tex} and \texttt{LIMITATIONS\_AND\_IMPACT.md})
    discusses:
    \begin{itemize}
      \item \textbf{Positive.} Lowering the barrier to RL post-training
            research; enabling reproducibility audits; surfacing
            platform-dependent variance that is typically hidden in
            single-platform papers.
      \item \textbf{Negative / risks.} RLHF reward hacking; alignment
            concerns when optimising proxy rewards; potential misuse of
            fine-tuned models; compute-access disparities that could
            limit replication to well-resourced groups.
      \item \textbf{Environmental.} Estimated 0.4--1.2\,kg
            CO$_2$-equivalent for the Modal H100 experiments; the Tinker
            API footprint is unestimable because the hardware mix is
            undisclosed.
    \end{itemize}}
    \end{enumerate}

% -- 11. Safeguards ------------------------------------------------------------
\item {\bf Safeguards}
    \begin{enumerate}
    \item Does the paper describe safeguards that have been put in place
          for responsible release of data or models that have been
          identified as requiring safeguards?
    \answerNA{The released artefacts are LoRA adapters fine-tuned from
    already-public base models (Qwen, Llama, Nemotron, GPT-OSS,
    Kimi-K2) on standard research datasets (GSM8K, HumanEval, synthetic
    tool-use).  No new high-risk capabilities are introduced relative to
    the base models; released checkpoints inherit the base models'
    licence terms and include standard usage disclaimers in the
    HuggingFace model cards
    (\texttt{huggingface/MODEL\_CARD\_TEMPLATE.md}).}
    \end{enumerate}

% -- 12. Licenses --------------------------------------------------------------
\item {\bf Licenses for existing assets}
    \begin{enumerate}
    \item Are the creators or original owners of assets (e.g., code,
          data, models) used in the paper properly credited, and are
          the licence and terms of use explicitly mentioned and properly
          respected?
    \answerYes{All third-party assets are credited with their licences:
    \begin{itemize}
      \item \textbf{GSM8K}~\cite{cobbe2021training}: MIT licence.
      \item \textbf{HumanEval} (OpenAI): MIT licence.
      \item \textbf{NoRobots} (HuggingFaceH4): CC-BY-NC-4.0 (used for
            the chat-SFT task only).
      \item \textbf{Synthetic tool-use data}: self-generated from public
            API documentation; no upstream licence restrictions.
      \item \textbf{Qwen model family}: Apache-2.0.
      \item \textbf{Llama model family}: Meta Llama Community Licence.
      \item \textbf{Nemotron, GPT-OSS, Kimi-K2}: used under their
            respective published licences; credits in
            Section~\ref{sec:setup}.
      \item \textbf{TRL, PEFT, Transformers}: Apache-2.0 (HuggingFace).
      \item \textbf{Modal}: commercial platform; terms of service
            respected.
      \item \textbf{Tinker API}: proprietary; used under standard API
            terms of service.
      \item \textbf{Our code}: \textbf{Apache-2.0} (see \texttt{LICENSE}
            in the repository root).
    \end{itemize}}
    \end{enumerate}

% -- 13. New Assets ------------------------------------------------------------
\item {\bf New Assets}
    \begin{enumerate}
    \item Are new assets introduced in the paper well documented and is
          the documentation provided alongside the assets?
    \answerYes{New assets and their documentation:
    \begin{itemize}
      \item The \ours{} benchmark suite (harness, evaluation scripts,
            analysis notebooks) at
            \url{https://github.com/arvindcr4/tinker-rl-lab}, released
            under Apache-2.0 and documented by \texttt{README.md},
            \texttt{REPRODUCE.md}, \texttt{ARTIFACT.md},
            \texttt{COMPUTE.md}, \texttt{BASELINES.md},
            \texttt{BENCHMARKS\_COMPARISON.md}, and
            \texttt{LIMITATIONS\_AND\_IMPACT.md}.
      \item LoRA adapter checkpoints from the successful Modal
            experiments at
            \url{https://huggingface.co/arvindcr4/tinker-rl-bench-*},
            each accompanied by a HuggingFace model card generated from
            \texttt{huggingface/MODEL\_CARD\_TEMPLATE.md}
            (intended use, training data, evaluation results, known
            limitations).
      \item All experiment logs on the public Weights~\&~Biases project
            \texttt{arvindcr4-pes-university/tinker-rl-lab-world-class}.
    \end{itemize}
    Checkpoints from JWT-interrupted Tinker runs are \emph{not} uploaded,
    because those jobs did not reach a valid final state
    (Section~\ref{sec:infra_failures}).}
    \end{enumerate}

% -- 14. Crowdsourcing ---------------------------------------------------------
\item {\bf Crowdsourcing and Research with Human Subjects}
    \begin{enumerate}
    \item For crowdsourcing experiments and research with human
          subjects, does the paper include the full text of instructions
          given to participants and screenshots, if applicable, as well
          as details about compensation?
    \answerNA{No crowdsourcing or human subjects research was
    conducted.}
    \end{enumerate}

% -- 15. IRB -------------------------------------------------------------------
\item {\bf Institutional Review Board (IRB) Approvals or Equivalent for
       Research with Human Subjects}
    \begin{enumerate}
    \item Does the paper describe potential risks incurred by study
          participants, whether compensation was provided to study
          participants, and whether informed consent was obtained?
    \answerNA{No human subjects research was conducted; IRB approval is
    not applicable.}
    \end{enumerate}

% -- 16. LLM usage -------------------------------------------------------------
\item {\bf Declaration of LLM usage}
    \begin{enumerate}
    \item Does the paper describe the usage of LLMs if it is an
          important, original, or non-standard component of the core
          methods?
    \answerYes{LLMs appear in this work in two distinct roles, both
    declared:
    \begin{itemize}
      \item \textbf{As subjects of evaluation.} The Qwen, Llama,
            Nemotron, GPT-OSS, and Kimi-K2 model families are the
            primary \emph{subjects} of the benchmark; their role as RL
            fine-tuning targets is described in
            Sections~\ref{sec:benchmark} and~\ref{sec:setup}.
      \item \textbf{As authoring/editing aids.} GitHub Copilot was used
            for boilerplate experiment scripts and LaTeX/BibTeX
            scaffolding.  ChatGPT Pro was used during revision to
            obtain reviewer-style critique of drafts; resulting
            suggestions were evaluated and, where accepted, manually
            incorporated.  All scientific claims, experimental design,
            numerical results, figures, and statistical analyses are
            human-authored and independently verified against the logged
            Weights~\&~Biases traces.
    \end{itemize}
    LLMs are not an \emph{original} or \emph{non-standard} component of
    the core methodology.}
    \end{enumerate}

\end{enumerate}

% (the file lives at paper/sections/checklist.tex so the relative path is
% correct when main.tex is compiled from paper/).
%
% Requirements already present in main.tex:
%   \newcommand{\answerPartial}[1][]{\textcolor{purple}{[Partial] #1}}
%   neurips_2026.sty provides \answerYes, \answerNo, \answerNA.
%
% The matching human-readable audit with page references lives at
%   NEURIPS_CHECKLIST_FINAL.md in the repository root.
% =============================================================================

\section*{NeurIPS Paper Checklist}

\begin{enumerate}

% -- 1. Claims -----------------------------------------------------------------
\item {\bf Claims}
    \begin{enumerate}
    \item Do the main claims made in the abstract and introduction
          accurately reflect the paper's contributions and scope?
    \answerYes{The abstract and Section~\ref{sec:intro} state five
    contributions: (i)~a 73-percentage-point implementation gap across
    7~libraries, (ii)~model-dependent GRPO/PPO preference, (iii)~frontier
    collapse on Nemotron-120B, (iv)~zero-variance fraction (ZVF) as a
    leading diagnostic of GRPO failure, and (v)~reward-trajectory
    proxies for KL-free monitoring.  Empirical support is given in
    Section~\ref{sec:results}: the cross-library gap in
    Section~\ref{sec:main_results}; the algorithm$\times$model
    interaction in Section~\ref{sec:ppo_grpo}; the frontier regime in
    Section~\ref{sec:frontier}; ZVF analysis in
    Section~\ref{sec:zvf-analysis}; and proxy metrics in
    Sections~\ref{sec:length_bias} and~\ref{sec:discuss_kl}.  Evidence is
    drawn from the 44~controlled experiments tabulated in
    Section~\ref{sec:results} and the 32-run Bitter Lesson extension of
    Section~\ref{sec:bitter_lesson}.  The scope is explicitly restricted
    to LoRA fine-tuning on three task families (math, code, tool-use)
    with evaluated scales of 0.6B--235B parameters
    (Section~\ref{sec:setup}); Section~\ref{sec:limitations} records
    that 11 of 14 Tinker runs were interrupted by JWT token expiry and
    4 of 6 Modal runs timed out, so reported numbers reflect the
    \emph{available} data and partial runs are marked $\dagger$.}
    \end{enumerate}

% -- 2. Limitations ------------------------------------------------------------
\item {\bf Limitations}
    \begin{enumerate}
    \item Does the paper discuss the limitations of the work performed by
          the authors?
    \answerYes{Section~\ref{sec:limitations} enumerates the limitations
    in detail:
    \begin{itemize}
      \item \textbf{Platform opacity.} Tinker API is a closed, serverless
            platform; GPU type, driver version, CUDA runtime, and
            scheduler are undisclosed, so Tinker-derived results cannot be
            independently reproduced without Tinker access
            (Section~\ref{sec:infra_failures}).
      \item \textbf{Infrastructure failures.} 11~of 14 Tinker runs were
            interrupted by JWT token expiry; 4~of 6 Modal runs hit
            timeouts or gradient-norm blow-ups.  Specific failed runs are
            named in Section~\ref{sec:infra_failures}.
      \item \textbf{Single-seed Tinker runs.} Cost constraints precluded
            multi-seed replication on Tinker; these results are treated
            as descriptive only (Section~\ref{sec:method_limits}).
      \item \textbf{LoRA-only evaluation.} Full fine-tuning and
            quantisation-aware training are not evaluated
            (Section~\ref{sec:method_limits}).
      \item \textbf{Train-set reward metric.} Primary Tinker metrics are
            training rewards, not held-out accuracy, with held-out
            evaluation only completed for the Modal Qwen3-32B / Qwen3.5-27B
            arms (Section~\ref{sec:method_limits}).
      \item \textbf{Proxy metrics in place of direct KL.}
            Section~\ref{sec:discuss_kl} discusses the limits of
            trajectory-based stability indices as substitutes for direct
            KL divergence.
    \end{itemize}}
    \end{enumerate}

% -- 3. Theory ------------------------------------------------------------------
\item {\bf Theory Assumptions and Proofs}
    \begin{enumerate}
    \item For each theoretical result, does the paper provide the full
          set of assumptions and a complete (or correct) proof?
    \answerNA{This is an empirical benchmarking paper; no formal theorems
    are claimed.  GRPO and PPO objectives stated in
    Section~\ref{sec:benchmark} are background definitions, not original
    theoretical results.}
    \end{enumerate}

% -- 4. Reproducibility ---------------------------------------------------------
\item {\bf Experimental Result Reproducibility}
    \begin{enumerate}
    \item Does the paper fully disclose all the information needed to
          reproduce the main experimental results of the paper to the
          extent that it affects the main claims and/or conclusions of
          the paper?
    \answerPartial{Reproducibility is \emph{partial by design} due to
    platform heterogeneity.
    \begin{itemize}
      \item \textbf{Modal experiments (fully reproducible).} Complete
            source code, a pinned \texttt{Dockerfile}, centralised seed
            management (\texttt{utils/seed.py}), and step-by-step
            commands are documented in \texttt{REPRODUCE.md} at
            \url{https://github.com/arvindcr4/tinker-rl-lab}.  Modal
            jobs run on explicitly provisioned NVIDIA H100~SXM5 (80\,GB)
            workers; TRL baselines run on NVIDIA L4 (24\,GB).  All
            Modal/TRL arms use 5~seeds ($s \in \{42, 123, 456, 789, 1024\}$)
            with mean $\pm$ SE and 95\,\% bootstrap CIs via
            \texttt{rliable}.
      \item \textbf{Tinker API experiments (not fully reproducible).}
            Tinker is closed-source with serverless GPU dispatch; GPU
            type, driver version, and scheduling policy are not exposed.
            We release our exact API scripts and configuration files, but
            independent replication requires a Tinker account and is
            subject to the same JWT expiry issues we observed
            (Section~\ref{sec:infra_failures}).  All Tinker-derived
            numbers in figures and tables are flagged with a ``$\dagger$''.
    \end{itemize}
    We claim \emph{Partial} rather than \emph{Yes}: the Tinker subset is
    irreducibly platform-dependent.}
    \end{enumerate}

% -- 5. Open access -------------------------------------------------------------
\item {\bf Open access to data and code}
    \begin{enumerate}
    \item Does the paper provide open access to the data and code, with
          sufficient instructions to faithfully reproduce the main
          experimental results, as described above?
    \answerYes{All code is publicly released under the
    \textbf{Apache License~2.0} at
    \url{https://github.com/arvindcr4/tinker-rl-lab} (mirrored at
    \url{https://github.com/pes-llm-research/tinker-rl-lab}).  LoRA
    adapter checkpoints from the successful Modal experiments are on
    HuggingFace Hub under
    \url{https://huggingface.co/arvindcr4/tinker-rl-bench-*} with model
    cards generated from \texttt{huggingface/MODEL\_CARD\_TEMPLATE.md}.
    All experiment runs are logged publicly at
    \url{https://wandb.ai/arvindcr4-pes-university/tinker-rl-lab-world-class};
    per-step reward, loss, gradient-norm, and (where available) KL
    metrics are visible without login.  Training datasets (GSM8K,
    HumanEval, NoRobots, synthetic tool-use) are public and cited in
    Table~\ref{tab:tasks}.  Tinker API credentials are \emph{not}
    included for security reasons; \texttt{README.md} explains how to
    obtain a Tinker API key.}
    \end{enumerate}

% -- 6. Experimental setting ---------------------------------------------------
\item {\bf Experimental Setting/Details}
    \begin{enumerate}
    \item Does the paper specify all the training and test details
          (e.g., data splits, hyperparameters, how they were chosen,
          type of optimizer) necessary to understand the results?
    \answerYes{Section~\ref{sec:setup} gives models, data splits, LoRA
    configuration (rank~32), optimiser (Adam, $\beta_1{=}0.9$,
    $\beta_2{=}0.95$, $\epsilon{=}10^{-8}$, learning rate $10^{-4}$),
    and the 5-seed Modal protocol.  Table~\ref{tab:hyperparams} maps
    these hyperparameters across all 7~libraries.
    Appendix~\ref{app:hyperparams} reports sweep ranges.  Per-task
    YAML configurations are version-controlled under
    \texttt{atropos/configs/}.  GSM8K uses the standard
    train/test partitions; HumanEval uses the full pass@1 suite.
    Tinker hyperparameters are released as API call scripts; the only
    unavailable information is the Tinker platform's internal hardware.}
    \end{enumerate}

% -- 7. Statistical significance -----------------------------------------------
\item {\bf Experiment Statistical Significance}
    \begin{enumerate}
    \item Does the paper report error bars suitably and correctly
          defined or other appropriate information about the statistical
          significance of the experiments?
    \answerPartial{For Modal/TRL arms with 5~seeds we report mean
    $\pm$ SE with 95\,\% bootstrap confidence intervals via
    \texttt{rliable}~\cite{agarwal2021deep}; a full power analysis
    (Table~\ref{tab:power-analysis}) and Benjamini--Hochberg-adjusted
    pairwise significance (Table~\ref{tab:bh-pvalues}) are reported in
    Section~\ref{sec:setup}.  The TRL-GRPO reference characterisation
    (Qwen2.5-0.5B, 5~seeds) reports $\bar{x}=0.734$,
    $\sigma=0.0703$, IQM $=0.747$, bootstrap 95\,\% CI $[0.672, 0.782]$.
    Tinker API experiments are single-seed (cost-constrained) and are
    explicitly labelled so in Sections~\ref{sec:main_results}
    and~\ref{sec:method_limits}; no statistical significance is claimed
    for Tinker-only comparisons.  We follow the recommendations of
    \citet{colas2019hitchhiker} and \citet{jordan2024benchmarking} on
    the limits of single-seed evaluation.  Because a meaningful subset
    of headline results is Tinker-only, the honest answer is
    \emph{Partial}.}
    \end{enumerate}

% -- 8. Compute resources ------------------------------------------------------
\item {\bf Experiments Compute Resources}
    \begin{enumerate}
    \item For each experiment, does the paper provide sufficient
          information on the computer resources (type of compute
          workers, memory, time of execution) needed to reproduce the
          experiments?
    \answerPartial{Section~4.4 and Appendix~\ref{app:compute} list GPU
    types and estimated GPU-hours for each experiment group; the
    complete per-experiment breakdown with wall-clock time and
    estimated cost is in \texttt{COMPUTE.md}.
    \begin{itemize}
      \item \textbf{Modal experiments.} Fully specified:
            H100~SXM5 (80\,GB) for recent runs and L4 (24\,GB) for the
            TRL reference sweep.  Wall-clock and GPU-hour estimates are
            reported.
      \item \textbf{Tinker API experiments.} Serverless dispatch masks
            the exact GPU type; GPU-hour figures are inferred from
            billing records and are flagged as approximate.
    \end{itemize}
    Total project cost is approximately 1{,}200 A100-equivalent GPU-hours
    across both platforms.  \emph{Partial} reflects the Tinker-side
    hardware opacity, which is a platform property we cannot resolve.}
    \end{enumerate}

% -- 9. Code of Ethics ---------------------------------------------------------
\item {\bf Code Of Ethics}
    \begin{enumerate}
    \item Does the research conducted in the paper conform, in every
          respect, with the NeurIPS Code of Ethics?
    \answerYes{The authors have reviewed the NeurIPS Code of Ethics and
    confirm compliance.  This work involves no human subjects, no
    private or sensitive data, and no models trained on proprietary
    corpora.  All training datasets are public research datasets
    (GSM8K, HumanEval, NoRobots, synthetic tool-use).  Broader societal
    impacts---including reward hacking, alignment failure modes of
    RLHF, and equity of compute access---are discussed in
    Section~\ref{sec:impact} and in \texttt{ethics\_statement.tex}.}
    \end{enumerate}

% -- 10. Broader Impacts -------------------------------------------------------
\item {\bf Broader Impacts}
    \begin{enumerate}
    \item Does the paper discuss both potential positive societal
          impacts and negative societal impacts of the work performed?
    \answerYes{Section~\ref{sec:impact} (plus
    \texttt{ethics\_statement.tex} and \texttt{LIMITATIONS\_AND\_IMPACT.md})
    discusses:
    \begin{itemize}
      \item \textbf{Positive.} Lowering the barrier to RL post-training
            research; enabling reproducibility audits; surfacing
            platform-dependent variance that is typically hidden in
            single-platform papers.
      \item \textbf{Negative / risks.} RLHF reward hacking; alignment
            concerns when optimising proxy rewards; potential misuse of
            fine-tuned models; compute-access disparities that could
            limit replication to well-resourced groups.
      \item \textbf{Environmental.} Estimated 0.4--1.2\,kg
            CO$_2$-equivalent for the Modal H100 experiments; the Tinker
            API footprint is unestimable because the hardware mix is
            undisclosed.
    \end{itemize}}
    \end{enumerate}

% -- 11. Safeguards ------------------------------------------------------------
\item {\bf Safeguards}
    \begin{enumerate}
    \item Does the paper describe safeguards that have been put in place
          for responsible release of data or models that have been
          identified as requiring safeguards?
    \answerNA{The released artefacts are LoRA adapters fine-tuned from
    already-public base models (Qwen, Llama, Nemotron, GPT-OSS,
    Kimi-K2) on standard research datasets (GSM8K, HumanEval, synthetic
    tool-use).  No new high-risk capabilities are introduced relative to
    the base models; released checkpoints inherit the base models'
    licence terms and include standard usage disclaimers in the
    HuggingFace model cards
    (\texttt{huggingface/MODEL\_CARD\_TEMPLATE.md}).}
    \end{enumerate}

% -- 12. Licenses --------------------------------------------------------------
\item {\bf Licenses for existing assets}
    \begin{enumerate}
    \item Are the creators or original owners of assets (e.g., code,
          data, models) used in the paper properly credited, and are
          the licence and terms of use explicitly mentioned and properly
          respected?
    \answerYes{All third-party assets are credited with their licences:
    \begin{itemize}
      \item \textbf{GSM8K}~\cite{cobbe2021training}: MIT licence.
      \item \textbf{HumanEval} (OpenAI): MIT licence.
      \item \textbf{NoRobots} (HuggingFaceH4): CC-BY-NC-4.0 (used for
            the chat-SFT task only).
      \item \textbf{Synthetic tool-use data}: self-generated from public
            API documentation; no upstream licence restrictions.
      \item \textbf{Qwen model family}: Apache-2.0.
      \item \textbf{Llama model family}: Meta Llama Community Licence.
      \item \textbf{Nemotron, GPT-OSS, Kimi-K2}: used under their
            respective published licences; credits in
            Section~\ref{sec:setup}.
      \item \textbf{TRL, PEFT, Transformers}: Apache-2.0 (HuggingFace).
      \item \textbf{Modal}: commercial platform; terms of service
            respected.
      \item \textbf{Tinker API}: proprietary; used under standard API
            terms of service.
      \item \textbf{Our code}: \textbf{Apache-2.0} (see \texttt{LICENSE}
            in the repository root).
    \end{itemize}}
    \end{enumerate}

% -- 13. New Assets ------------------------------------------------------------
\item {\bf New Assets}
    \begin{enumerate}
    \item Are new assets introduced in the paper well documented and is
          the documentation provided alongside the assets?
    \answerYes{New assets and their documentation:
    \begin{itemize}
      \item The \ours{} benchmark suite (harness, evaluation scripts,
            analysis notebooks) at
            \url{https://github.com/arvindcr4/tinker-rl-lab}, released
            under Apache-2.0 and documented by \texttt{README.md},
            \texttt{REPRODUCE.md}, \texttt{ARTIFACT.md},
            \texttt{COMPUTE.md}, \texttt{BASELINES.md},
            \texttt{BENCHMARKS\_COMPARISON.md}, and
            \texttt{LIMITATIONS\_AND\_IMPACT.md}.
      \item LoRA adapter checkpoints from the successful Modal
            experiments at
            \url{https://huggingface.co/arvindcr4/tinker-rl-bench-*},
            each accompanied by a HuggingFace model card generated from
            \texttt{huggingface/MODEL\_CARD\_TEMPLATE.md}
            (intended use, training data, evaluation results, known
            limitations).
      \item All experiment logs on the public Weights~\&~Biases project
            \texttt{arvindcr4-pes-university/tinker-rl-lab-world-class}.
    \end{itemize}
    Checkpoints from JWT-interrupted Tinker runs are \emph{not} uploaded,
    because those jobs did not reach a valid final state
    (Section~\ref{sec:infra_failures}).}
    \end{enumerate}

% -- 14. Crowdsourcing ---------------------------------------------------------
\item {\bf Crowdsourcing and Research with Human Subjects}
    \begin{enumerate}
    \item For crowdsourcing experiments and research with human
          subjects, does the paper include the full text of instructions
          given to participants and screenshots, if applicable, as well
          as details about compensation?
    \answerNA{No crowdsourcing or human subjects research was
    conducted.}
    \end{enumerate}

% -- 15. IRB -------------------------------------------------------------------
\item {\bf Institutional Review Board (IRB) Approvals or Equivalent for
       Research with Human Subjects}
    \begin{enumerate}
    \item Does the paper describe potential risks incurred by study
          participants, whether compensation was provided to study
          participants, and whether informed consent was obtained?
    \answerNA{No human subjects research was conducted; IRB approval is
    not applicable.}
    \end{enumerate}

% -- 16. LLM usage -------------------------------------------------------------
\item {\bf Declaration of LLM usage}
    \begin{enumerate}
    \item Does the paper describe the usage of LLMs if it is an
          important, original, or non-standard component of the core
          methods?
    \answerYes{LLMs appear in this work in two distinct roles, both
    declared:
    \begin{itemize}
      \item \textbf{As subjects of evaluation.} The Qwen, Llama,
            Nemotron, GPT-OSS, and Kimi-K2 model families are the
            primary \emph{subjects} of the benchmark; their role as RL
            fine-tuning targets is described in
            Sections~\ref{sec:benchmark} and~\ref{sec:setup}.
      \item \textbf{As authoring/editing aids.} GitHub Copilot was used
            for boilerplate experiment scripts and LaTeX/BibTeX
            scaffolding.  ChatGPT Pro was used during revision to
            obtain reviewer-style critique of drafts; resulting
            suggestions were evaluated and, where accepted, manually
            incorporated.  All scientific claims, experimental design,
            numerical results, figures, and statistical analyses are
            human-authored and independently verified against the logged
            Weights~\&~Biases traces.
    \end{itemize}
    LLMs are not an \emph{original} or \emph{non-standard} component of
    the core methodology.}
    \end{enumerate}

\end{enumerate}

% (the file lives at paper/sections/checklist.tex so the relative path is
% correct when main.tex is compiled from paper/).
%
% Requirements already present in main.tex:
%   \newcommand{\answerPartial}[1][]{\textcolor{purple}{[Partial] #1}}
%   neurips_2026.sty provides \answerYes, \answerNo, \answerNA.
%
% The matching human-readable audit with page references lives at
%   NEURIPS_CHECKLIST_FINAL.md in the repository root.
% =============================================================================

\section*{NeurIPS Paper Checklist}

\begin{enumerate}

% -- 1. Claims -----------------------------------------------------------------
\item {\bf Claims}
    \begin{enumerate}
    \item Do the main claims made in the abstract and introduction
          accurately reflect the paper's contributions and scope?
    \answerYes{The abstract and Section~\ref{sec:intro} state five
    contributions: (i)~a 73-percentage-point implementation gap across
    7~libraries, (ii)~model-dependent GRPO/PPO preference, (iii)~frontier
    collapse on Nemotron-120B, (iv)~zero-variance fraction (ZVF) as a
    leading diagnostic of GRPO failure, and (v)~reward-trajectory
    proxies for KL-free monitoring.  Empirical support is given in
    Section~\ref{sec:results}: the cross-library gap in
    Section~\ref{sec:main_results}; the algorithm$\times$model
    interaction in Section~\ref{sec:ppo_grpo}; the frontier regime in
    Section~\ref{sec:frontier}; ZVF analysis in
    Section~\ref{sec:zvf-analysis}; and proxy metrics in
    Sections~\ref{sec:length_bias} and~\ref{sec:discuss_kl}.  Evidence is
    drawn from the 44~controlled experiments tabulated in
    Section~\ref{sec:results} and the 32-run Bitter Lesson extension of
    Section~\ref{sec:bitter_lesson}.  The scope is explicitly restricted
    to LoRA fine-tuning on three task families (math, code, tool-use)
    with evaluated scales of 0.6B--235B parameters
    (Section~\ref{sec:setup}); Section~\ref{sec:limitations} records
    that 11 of 14 Tinker runs were interrupted by JWT token expiry and
    4 of 6 Modal runs timed out, so reported numbers reflect the
    \emph{available} data and partial runs are marked $\dagger$.}
    \end{enumerate}

% -- 2. Limitations ------------------------------------------------------------
\item {\bf Limitations}
    \begin{enumerate}
    \item Does the paper discuss the limitations of the work performed by
          the authors?
    \answerYes{Section~\ref{sec:limitations} enumerates the limitations
    in detail:
    \begin{itemize}
      \item \textbf{Platform opacity.} Tinker API is a closed, serverless
            platform; GPU type, driver version, CUDA runtime, and
            scheduler are undisclosed, so Tinker-derived results cannot be
            independently reproduced without Tinker access
            (Section~\ref{sec:infra_failures}).
      \item \textbf{Infrastructure failures.} 11~of 14 Tinker runs were
            interrupted by JWT token expiry; 4~of 6 Modal runs hit
            timeouts or gradient-norm blow-ups.  Specific failed runs are
            named in Section~\ref{sec:infra_failures}.
      \item \textbf{Single-seed Tinker runs.} Cost constraints precluded
            multi-seed replication on Tinker; these results are treated
            as descriptive only (Section~\ref{sec:method_limits}).
      \item \textbf{LoRA-only evaluation.} Full fine-tuning and
            quantisation-aware training are not evaluated
            (Section~\ref{sec:method_limits}).
      \item \textbf{Train-set reward metric.} Primary Tinker metrics are
            training rewards, not held-out accuracy, with held-out
            evaluation only completed for the Modal Qwen3-32B / Qwen3.5-27B
            arms (Section~\ref{sec:method_limits}).
      \item \textbf{Proxy metrics in place of direct KL.}
            Section~\ref{sec:discuss_kl} discusses the limits of
            trajectory-based stability indices as substitutes for direct
            KL divergence.
    \end{itemize}}
    \end{enumerate}

% -- 3. Theory ------------------------------------------------------------------
\item {\bf Theory Assumptions and Proofs}
    \begin{enumerate}
    \item For each theoretical result, does the paper provide the full
          set of assumptions and a complete (or correct) proof?
    \answerNA{This is an empirical benchmarking paper; no formal theorems
    are claimed.  GRPO and PPO objectives stated in
    Section~\ref{sec:benchmark} are background definitions, not original
    theoretical results.}
    \end{enumerate}

% -- 4. Reproducibility ---------------------------------------------------------
\item {\bf Experimental Result Reproducibility}
    \begin{enumerate}
    \item Does the paper fully disclose all the information needed to
          reproduce the main experimental results of the paper to the
          extent that it affects the main claims and/or conclusions of
          the paper?
    \answerPartial{Reproducibility is \emph{partial by design} due to
    platform heterogeneity.
    \begin{itemize}
      \item \textbf{Modal experiments (fully reproducible).} Complete
            source code, a pinned \texttt{Dockerfile}, centralised seed
            management (\texttt{utils/seed.py}), and step-by-step
            commands are documented in \texttt{REPRODUCE.md} at
            \url{https://github.com/arvindcr4/tinker-rl-lab}.  Modal
            jobs run on explicitly provisioned NVIDIA H100~SXM5 (80\,GB)
            workers; TRL baselines run on NVIDIA L4 (24\,GB).  All
            Modal/TRL arms use 5~seeds ($s \in \{42, 123, 456, 789, 1024\}$)
            with mean $\pm$ SE and 95\,\% bootstrap CIs via
            \texttt{rliable}.
      \item \textbf{Tinker API experiments (not fully reproducible).}
            Tinker is closed-source with serverless GPU dispatch; GPU
            type, driver version, and scheduling policy are not exposed.
            We release our exact API scripts and configuration files, but
            independent replication requires a Tinker account and is
            subject to the same JWT expiry issues we observed
            (Section~\ref{sec:infra_failures}).  All Tinker-derived
            numbers in figures and tables are flagged with a ``$\dagger$''.
    \end{itemize}
    We claim \emph{Partial} rather than \emph{Yes}: the Tinker subset is
    irreducibly platform-dependent.}
    \end{enumerate}

% -- 5. Open access -------------------------------------------------------------
\item {\bf Open access to data and code}
    \begin{enumerate}
    \item Does the paper provide open access to the data and code, with
          sufficient instructions to faithfully reproduce the main
          experimental results, as described above?
    \answerYes{All code is publicly released under the
    \textbf{Apache License~2.0} at
    \url{https://github.com/arvindcr4/tinker-rl-lab} (mirrored at
    \url{https://github.com/pes-llm-research/tinker-rl-lab}).  LoRA
    adapter checkpoints from the successful Modal experiments are on
    HuggingFace Hub under
    \url{https://huggingface.co/arvindcr4/tinker-rl-bench-*} with model
    cards generated from \texttt{huggingface/MODEL\_CARD\_TEMPLATE.md}.
    All experiment runs are logged publicly at
    \url{https://wandb.ai/arvindcr4-pes-university/tinker-rl-lab-world-class};
    per-step reward, loss, gradient-norm, and (where available) KL
    metrics are visible without login.  Training datasets (GSM8K,
    HumanEval, NoRobots, synthetic tool-use) are public and cited in
    Table~\ref{tab:tasks}.  Tinker API credentials are \emph{not}
    included for security reasons; \texttt{README.md} explains how to
    obtain a Tinker API key.}
    \end{enumerate}

% -- 6. Experimental setting ---------------------------------------------------
\item {\bf Experimental Setting/Details}
    \begin{enumerate}
    \item Does the paper specify all the training and test details
          (e.g., data splits, hyperparameters, how they were chosen,
          type of optimizer) necessary to understand the results?
    \answerYes{Section~\ref{sec:setup} gives models, data splits, LoRA
    configuration (rank~32), optimiser (Adam, $\beta_1{=}0.9$,
    $\beta_2{=}0.95$, $\epsilon{=}10^{-8}$, learning rate $10^{-4}$),
    and the 5-seed Modal protocol.  Table~\ref{tab:hyperparams} maps
    these hyperparameters across all 7~libraries.
    Appendix~\ref{app:hyperparams} reports sweep ranges.  Per-task
    YAML configurations are version-controlled under
    \texttt{atropos/configs/}.  GSM8K uses the standard
    train/test partitions; HumanEval uses the full pass@1 suite.
    Tinker hyperparameters are released as API call scripts; the only
    unavailable information is the Tinker platform's internal hardware.}
    \end{enumerate}

% -- 7. Statistical significance -----------------------------------------------
\item {\bf Experiment Statistical Significance}
    \begin{enumerate}
    \item Does the paper report error bars suitably and correctly
          defined or other appropriate information about the statistical
          significance of the experiments?
    \answerPartial{For Modal/TRL arms with 5~seeds we report mean
    $\pm$ SE with 95\,\% bootstrap confidence intervals via
    \texttt{rliable}~\cite{agarwal2021deep}; a full power analysis
    (Table~\ref{tab:power-analysis}) and Benjamini--Hochberg-adjusted
    pairwise significance (Table~\ref{tab:bh-pvalues}) are reported in
    Section~\ref{sec:setup}.  The TRL-GRPO reference characterisation
    (Qwen2.5-0.5B, 5~seeds) reports $\bar{x}=0.734$,
    $\sigma=0.0703$, IQM $=0.747$, bootstrap 95\,\% CI $[0.672, 0.782]$.
    Tinker API experiments are single-seed (cost-constrained) and are
    explicitly labelled so in Sections~\ref{sec:main_results}
    and~\ref{sec:method_limits}; no statistical significance is claimed
    for Tinker-only comparisons.  We follow the recommendations of
    \citet{colas2019hitchhiker} and \citet{jordan2024benchmarking} on
    the limits of single-seed evaluation.  Because a meaningful subset
    of headline results is Tinker-only, the honest answer is
    \emph{Partial}.}
    \end{enumerate}

% -- 8. Compute resources ------------------------------------------------------
\item {\bf Experiments Compute Resources}
    \begin{enumerate}
    \item For each experiment, does the paper provide sufficient
          information on the computer resources (type of compute
          workers, memory, time of execution) needed to reproduce the
          experiments?
    \answerPartial{Section~4.4 and Appendix~\ref{app:compute} list GPU
    types and estimated GPU-hours for each experiment group; the
    complete per-experiment breakdown with wall-clock time and
    estimated cost is in \texttt{COMPUTE.md}.
    \begin{itemize}
      \item \textbf{Modal experiments.} Fully specified:
            H100~SXM5 (80\,GB) for recent runs and L4 (24\,GB) for the
            TRL reference sweep.  Wall-clock and GPU-hour estimates are
            reported.
      \item \textbf{Tinker API experiments.} Serverless dispatch masks
            the exact GPU type; GPU-hour figures are inferred from
            billing records and are flagged as approximate.
    \end{itemize}
    Total project cost is approximately 1{,}200 A100-equivalent GPU-hours
    across both platforms.  \emph{Partial} reflects the Tinker-side
    hardware opacity, which is a platform property we cannot resolve.}
    \end{enumerate}

% -- 9. Code of Ethics ---------------------------------------------------------
\item {\bf Code Of Ethics}
    \begin{enumerate}
    \item Does the research conducted in the paper conform, in every
          respect, with the NeurIPS Code of Ethics?
    \answerYes{The authors have reviewed the NeurIPS Code of Ethics and
    confirm compliance.  This work involves no human subjects, no
    private or sensitive data, and no models trained on proprietary
    corpora.  All training datasets are public research datasets
    (GSM8K, HumanEval, NoRobots, synthetic tool-use).  Broader societal
    impacts---including reward hacking, alignment failure modes of
    RLHF, and equity of compute access---are discussed in
    Section~\ref{sec:impact} and in \texttt{ethics\_statement.tex}.}
    \end{enumerate}

% -- 10. Broader Impacts -------------------------------------------------------
\item {\bf Broader Impacts}
    \begin{enumerate}
    \item Does the paper discuss both potential positive societal
          impacts and negative societal impacts of the work performed?
    \answerYes{Section~\ref{sec:impact} (plus
    \texttt{ethics\_statement.tex} and \texttt{LIMITATIONS\_AND\_IMPACT.md})
    discusses:
    \begin{itemize}
      \item \textbf{Positive.} Lowering the barrier to RL post-training
            research; enabling reproducibility audits; surfacing
            platform-dependent variance that is typically hidden in
            single-platform papers.
      \item \textbf{Negative / risks.} RLHF reward hacking; alignment
            concerns when optimising proxy rewards; potential misuse of
            fine-tuned models; compute-access disparities that could
            limit replication to well-resourced groups.
      \item \textbf{Environmental.} Estimated 0.4--1.2\,kg
            CO$_2$-equivalent for the Modal H100 experiments; the Tinker
            API footprint is unestimable because the hardware mix is
            undisclosed.
    \end{itemize}}
    \end{enumerate}

% -- 11. Safeguards ------------------------------------------------------------
\item {\bf Safeguards}
    \begin{enumerate}
    \item Does the paper describe safeguards that have been put in place
          for responsible release of data or models that have been
          identified as requiring safeguards?
    \answerNA{The released artefacts are LoRA adapters fine-tuned from
    already-public base models (Qwen, Llama, Nemotron, GPT-OSS,
    Kimi-K2) on standard research datasets (GSM8K, HumanEval, synthetic
    tool-use).  No new high-risk capabilities are introduced relative to
    the base models; released checkpoints inherit the base models'
    licence terms and include standard usage disclaimers in the
    HuggingFace model cards
    (\texttt{huggingface/MODEL\_CARD\_TEMPLATE.md}).}
    \end{enumerate}

% -- 12. Licenses --------------------------------------------------------------
\item {\bf Licenses for existing assets}
    \begin{enumerate}
    \item Are the creators or original owners of assets (e.g., code,
          data, models) used in the paper properly credited, and are
          the licence and terms of use explicitly mentioned and properly
          respected?
    \answerYes{All third-party assets are credited with their licences:
    \begin{itemize}
      \item \textbf{GSM8K}~\cite{cobbe2021training}: MIT licence.
      \item \textbf{HumanEval} (OpenAI): MIT licence.
      \item \textbf{NoRobots} (HuggingFaceH4): CC-BY-NC-4.0 (used for
            the chat-SFT task only).
      \item \textbf{Synthetic tool-use data}: self-generated from public
            API documentation; no upstream licence restrictions.
      \item \textbf{Qwen model family}: Apache-2.0.
      \item \textbf{Llama model family}: Meta Llama Community Licence.
      \item \textbf{Nemotron, GPT-OSS, Kimi-K2}: used under their
            respective published licences; credits in
            Section~\ref{sec:setup}.
      \item \textbf{TRL, PEFT, Transformers}: Apache-2.0 (HuggingFace).
      \item \textbf{Modal}: commercial platform; terms of service
            respected.
      \item \textbf{Tinker API}: proprietary; used under standard API
            terms of service.
      \item \textbf{Our code}: \textbf{Apache-2.0} (see \texttt{LICENSE}
            in the repository root).
    \end{itemize}}
    \end{enumerate}

% -- 13. New Assets ------------------------------------------------------------
\item {\bf New Assets}
    \begin{enumerate}
    \item Are new assets introduced in the paper well documented and is
          the documentation provided alongside the assets?
    \answerYes{New assets and their documentation:
    \begin{itemize}
      \item The \ours{} benchmark suite (harness, evaluation scripts,
            analysis notebooks) at
            \url{https://github.com/arvindcr4/tinker-rl-lab}, released
            under Apache-2.0 and documented by \texttt{README.md},
            \texttt{REPRODUCE.md}, \texttt{ARTIFACT.md},
            \texttt{COMPUTE.md}, \texttt{BASELINES.md},
            \texttt{BENCHMARKS\_COMPARISON.md}, and
            \texttt{LIMITATIONS\_AND\_IMPACT.md}.
      \item LoRA adapter checkpoints from the successful Modal
            experiments at
            \url{https://huggingface.co/arvindcr4/tinker-rl-bench-*},
            each accompanied by a HuggingFace model card generated from
            \texttt{huggingface/MODEL\_CARD\_TEMPLATE.md}
            (intended use, training data, evaluation results, known
            limitations).
      \item All experiment logs on the public Weights~\&~Biases project
            \texttt{arvindcr4-pes-university/tinker-rl-lab-world-class}.
    \end{itemize}
    Checkpoints from JWT-interrupted Tinker runs are \emph{not} uploaded,
    because those jobs did not reach a valid final state
    (Section~\ref{sec:infra_failures}).}
    \end{enumerate}

% -- 14. Crowdsourcing ---------------------------------------------------------
\item {\bf Crowdsourcing and Research with Human Subjects}
    \begin{enumerate}
    \item For crowdsourcing experiments and research with human
          subjects, does the paper include the full text of instructions
          given to participants and screenshots, if applicable, as well
          as details about compensation?
    \answerNA{No crowdsourcing or human subjects research was
    conducted.}
    \end{enumerate}

% -- 15. IRB -------------------------------------------------------------------
\item {\bf Institutional Review Board (IRB) Approvals or Equivalent for
       Research with Human Subjects}
    \begin{enumerate}
    \item Does the paper describe potential risks incurred by study
          participants, whether compensation was provided to study
          participants, and whether informed consent was obtained?
    \answerNA{No human subjects research was conducted; IRB approval is
    not applicable.}
    \end{enumerate}

% -- 16. LLM usage -------------------------------------------------------------
\item {\bf Declaration of LLM usage}
    \begin{enumerate}
    \item Does the paper describe the usage of LLMs if it is an
          important, original, or non-standard component of the core
          methods?
    \answerYes{LLMs appear in this work in two distinct roles, both
    declared:
    \begin{itemize}
      \item \textbf{As subjects of evaluation.} The Qwen, Llama,
            Nemotron, GPT-OSS, and Kimi-K2 model families are the
            primary \emph{subjects} of the benchmark; their role as RL
            fine-tuning targets is described in
            Sections~\ref{sec:benchmark} and~\ref{sec:setup}.
      \item \textbf{As authoring/editing aids.} GitHub Copilot was used
            for boilerplate experiment scripts and LaTeX/BibTeX
            scaffolding.  ChatGPT Pro was used during revision to
            obtain reviewer-style critique of drafts; resulting
            suggestions were evaluated and, where accepted, manually
            incorporated.  All scientific claims, experimental design,
            numerical results, figures, and statistical analyses are
            human-authored and independently verified against the logged
            Weights~\&~Biases traces.
    \end{itemize}
    LLMs are not an \emph{original} or \emph{non-standard} component of
    the core methodology.}
    \end{enumerate}

\end{enumerate}

% (the file lives at paper/sections/checklist.tex so the relative path is
% correct when main.tex is compiled from paper/).
%
% Requirements already present in main.tex:
%   \newcommand{\answerPartial}[1][]{\textcolor{purple}{[Partial] #1}}
%   neurips_2026.sty provides \answerYes, \answerNo, \answerNA.
%
% The matching human-readable audit with page references lives at
%   NEURIPS_CHECKLIST_FINAL.md in the repository root.
% =============================================================================

\section*{NeurIPS Paper Checklist}

\begin{enumerate}

% -- 1. Claims -----------------------------------------------------------------
\item {\bf Claims}
    \begin{enumerate}
    \item Do the main claims made in the abstract and introduction
          accurately reflect the paper's contributions and scope?
    \answerYes{The abstract and Section~\ref{sec:intro} state five
    contributions: (i)~a 73-percentage-point implementation gap across
    7~libraries, (ii)~model-dependent GRPO/PPO preference, (iii)~frontier
    collapse on Nemotron-120B, (iv)~zero-variance fraction (ZVF) as a
    leading diagnostic of GRPO failure, and (v)~reward-trajectory
    proxies for KL-free monitoring.  Empirical support is given in
    Section~\ref{sec:results}: the cross-library gap in
    Section~\ref{sec:main_results}; the algorithm$\times$model
    interaction in Section~\ref{sec:ppo_grpo}; the frontier regime in
    Section~\ref{sec:frontier}; ZVF analysis in
    Section~\ref{sec:zvf-analysis}; and proxy metrics in
    Sections~\ref{sec:length_bias} and~\ref{sec:discuss_kl}.  Evidence is
    drawn from the 44~controlled experiments tabulated in
    Section~\ref{sec:results} and the 32-run Bitter Lesson extension of
    Section~\ref{sec:bitter_lesson}.  The scope is explicitly restricted
    to LoRA fine-tuning on three task families (math, code, tool-use)
    with evaluated scales of 0.6B--235B parameters
    (Section~\ref{sec:setup}); Section~\ref{sec:limitations} records
    that 11 of 14 Tinker runs were interrupted by JWT token expiry and
    4 of 6 Modal runs timed out, so reported numbers reflect the
    \emph{available} data and partial runs are marked $\dagger$.}
    \end{enumerate}

% -- 2. Limitations ------------------------------------------------------------
\item {\bf Limitations}
    \begin{enumerate}
    \item Does the paper discuss the limitations of the work performed by
          the authors?
    \answerYes{Section~\ref{sec:limitations} enumerates the limitations
    in detail:
    \begin{itemize}
      \item \textbf{Platform opacity.} Tinker API is a closed, serverless
            platform; GPU type, driver version, CUDA runtime, and
            scheduler are undisclosed, so Tinker-derived results cannot be
            independently reproduced without Tinker access
            (Section~\ref{sec:infra_failures}).
      \item \textbf{Infrastructure failures.} 11~of 14 Tinker runs were
            interrupted by JWT token expiry; 4~of 6 Modal runs hit
            timeouts or gradient-norm blow-ups.  Specific failed runs are
            named in Section~\ref{sec:infra_failures}.
      \item \textbf{Single-seed Tinker runs.} Cost constraints precluded
            multi-seed replication on Tinker; these results are treated
            as descriptive only (Section~\ref{sec:method_limits}).
      \item \textbf{LoRA-only evaluation.} Full fine-tuning and
            quantisation-aware training are not evaluated
            (Section~\ref{sec:method_limits}).
      \item \textbf{Train-set reward metric.} Primary Tinker metrics are
            training rewards, not held-out accuracy, with held-out
            evaluation only completed for the Modal Qwen3-32B / Qwen3.5-27B
            arms (Section~\ref{sec:method_limits}).
      \item \textbf{Proxy metrics in place of direct KL.}
            Section~\ref{sec:discuss_kl} discusses the limits of
            trajectory-based stability indices as substitutes for direct
            KL divergence.
    \end{itemize}}
    \end{enumerate}

% -- 3. Theory ------------------------------------------------------------------
\item {\bf Theory Assumptions and Proofs}
    \begin{enumerate}
    \item For each theoretical result, does the paper provide the full
          set of assumptions and a complete (or correct) proof?
    \answerNA{This is an empirical benchmarking paper; no formal theorems
    are claimed.  GRPO and PPO objectives stated in
    Section~\ref{sec:benchmark} are background definitions, not original
    theoretical results.}
    \end{enumerate}

% -- 4. Reproducibility ---------------------------------------------------------
\item {\bf Experimental Result Reproducibility}
    \begin{enumerate}
    \item Does the paper fully disclose all the information needed to
          reproduce the main experimental results of the paper to the
          extent that it affects the main claims and/or conclusions of
          the paper?
    \answerPartial{Reproducibility is \emph{partial by design} due to
    platform heterogeneity.
    \begin{itemize}
      \item \textbf{Modal experiments (fully reproducible).} Complete
            source code, a pinned \texttt{Dockerfile}, centralised seed
            management (\texttt{utils/seed.py}), and step-by-step
            commands are documented in \texttt{REPRODUCE.md} at
            \url{https://github.com/arvindcr4/tinker-rl-lab}.  Modal
            jobs run on explicitly provisioned NVIDIA H100~SXM5 (80\,GB)
            workers; TRL baselines run on NVIDIA L4 (24\,GB).  All
            Modal/TRL arms use 5~seeds ($s \in \{42, 123, 456, 789, 1024\}$)
            with mean $\pm$ SE and 95\,\% bootstrap CIs via
            \texttt{rliable}.
      \item \textbf{Tinker API experiments (not fully reproducible).}
            Tinker is closed-source with serverless GPU dispatch; GPU
            type, driver version, and scheduling policy are not exposed.
            We release our exact API scripts and configuration files, but
            independent replication requires a Tinker account and is
            subject to the same JWT expiry issues we observed
            (Section~\ref{sec:infra_failures}).  All Tinker-derived
            numbers in figures and tables are flagged with a ``$\dagger$''.
    \end{itemize}
    We claim \emph{Partial} rather than \emph{Yes}: the Tinker subset is
    irreducibly platform-dependent.}
    \end{enumerate}

% -- 5. Open access -------------------------------------------------------------
\item {\bf Open access to data and code}
    \begin{enumerate}
    \item Does the paper provide open access to the data and code, with
          sufficient instructions to faithfully reproduce the main
          experimental results, as described above?
    \answerYes{All code is publicly released under the
    \textbf{Apache License~2.0} at
    \url{https://github.com/arvindcr4/tinker-rl-lab} (mirrored at
    \url{https://github.com/pes-llm-research/tinker-rl-lab}).  LoRA
    adapter checkpoints from the successful Modal experiments are on
    HuggingFace Hub under
    \url{https://huggingface.co/arvindcr4/tinker-rl-bench-*} with model
    cards generated from \texttt{huggingface/MODEL\_CARD\_TEMPLATE.md}.
    All experiment runs are logged publicly at
    \url{https://wandb.ai/arvindcr4-pes-university/tinker-rl-lab-world-class};
    per-step reward, loss, gradient-norm, and (where available) KL
    metrics are visible without login.  Training datasets (GSM8K,
    HumanEval, NoRobots, synthetic tool-use) are public and cited in
    Table~\ref{tab:tasks}.  Tinker API credentials are \emph{not}
    included for security reasons; \texttt{README.md} explains how to
    obtain a Tinker API key.}
    \end{enumerate}

% -- 6. Experimental setting ---------------------------------------------------
\item {\bf Experimental Setting/Details}
    \begin{enumerate}
    \item Does the paper specify all the training and test details
          (e.g., data splits, hyperparameters, how they were chosen,
          type of optimizer) necessary to understand the results?
    \answerYes{Section~\ref{sec:setup} gives models, data splits, LoRA
    configuration (rank~32), optimiser (Adam, $\beta_1{=}0.9$,
    $\beta_2{=}0.95$, $\epsilon{=}10^{-8}$, learning rate $10^{-4}$),
    and the 5-seed Modal protocol.  Table~\ref{tab:hyperparams} maps
    these hyperparameters across all 7~libraries.
    Appendix~\ref{app:hyperparams} reports sweep ranges.  Per-task
    YAML configurations are version-controlled under
    \texttt{atropos/configs/}.  GSM8K uses the standard
    train/test partitions; HumanEval uses the full pass@1 suite.
    Tinker hyperparameters are released as API call scripts; the only
    unavailable information is the Tinker platform's internal hardware.}
    \end{enumerate}

% -- 7. Statistical significance -----------------------------------------------
\item {\bf Experiment Statistical Significance}
    \begin{enumerate}
    \item Does the paper report error bars suitably and correctly
          defined or other appropriate information about the statistical
          significance of the experiments?
    \answerPartial{For Modal/TRL arms with 5~seeds we report mean
    $\pm$ SE with 95\,\% bootstrap confidence intervals via
    \texttt{rliable}~\cite{agarwal2021deep}; a full power analysis
    (Table~\ref{tab:power-analysis}) and Benjamini--Hochberg-adjusted
    pairwise significance (Table~\ref{tab:bh-pvalues}) are reported in
    Section~\ref{sec:setup}.  The TRL-GRPO reference characterisation
    (Qwen2.5-0.5B, 5~seeds) reports $\bar{x}=0.734$,
    $\sigma=0.0703$, IQM $=0.747$, bootstrap 95\,\% CI $[0.672, 0.782]$.
    Tinker API experiments are single-seed (cost-constrained) and are
    explicitly labelled so in Sections~\ref{sec:main_results}
    and~\ref{sec:method_limits}; no statistical significance is claimed
    for Tinker-only comparisons.  We follow the recommendations of
    \citet{colas2019hitchhiker} and \citet{jordan2024benchmarking} on
    the limits of single-seed evaluation.  Because a meaningful subset
    of headline results is Tinker-only, the honest answer is
    \emph{Partial}.}
    \end{enumerate}

% -- 8. Compute resources ------------------------------------------------------
\item {\bf Experiments Compute Resources}
    \begin{enumerate}
    \item For each experiment, does the paper provide sufficient
          information on the computer resources (type of compute
          workers, memory, time of execution) needed to reproduce the
          experiments?
    \answerPartial{Section~4.4 and Appendix~\ref{app:compute} list GPU
    types and estimated GPU-hours for each experiment group; the
    complete per-experiment breakdown with wall-clock time and
    estimated cost is in \texttt{COMPUTE.md}.
    \begin{itemize}
      \item \textbf{Modal experiments.} Fully specified:
            H100~SXM5 (80\,GB) for recent runs and L4 (24\,GB) for the
            TRL reference sweep.  Wall-clock and GPU-hour estimates are
            reported.
      \item \textbf{Tinker API experiments.} Serverless dispatch masks
            the exact GPU type; GPU-hour figures are inferred from
            billing records and are flagged as approximate.
    \end{itemize}
    Total project cost is approximately 1{,}200 A100-equivalent GPU-hours
    across both platforms.  \emph{Partial} reflects the Tinker-side
    hardware opacity, which is a platform property we cannot resolve.}
    \end{enumerate}

% -- 9. Code of Ethics ---------------------------------------------------------
\item {\bf Code Of Ethics}
    \begin{enumerate}
    \item Does the research conducted in the paper conform, in every
          respect, with the NeurIPS Code of Ethics?
    \answerYes{The authors have reviewed the NeurIPS Code of Ethics and
    confirm compliance.  This work involves no human subjects, no
    private or sensitive data, and no models trained on proprietary
    corpora.  All training datasets are public research datasets
    (GSM8K, HumanEval, NoRobots, synthetic tool-use).  Broader societal
    impacts---including reward hacking, alignment failure modes of
    RLHF, and equity of compute access---are discussed in
    Section~\ref{sec:impact} and in \texttt{ethics\_statement.tex}.}
    \end{enumerate}

% -- 10. Broader Impacts -------------------------------------------------------
\item {\bf Broader Impacts}
    \begin{enumerate}
    \item Does the paper discuss both potential positive societal
          impacts and negative societal impacts of the work performed?
    \answerYes{Section~\ref{sec:impact} (plus
    \texttt{ethics\_statement.tex} and \texttt{LIMITATIONS\_AND\_IMPACT.md})
    discusses:
    \begin{itemize}
      \item \textbf{Positive.} Lowering the barrier to RL post-training
            research; enabling reproducibility audits; surfacing
            platform-dependent variance that is typically hidden in
            single-platform papers.
      \item \textbf{Negative / risks.} RLHF reward hacking; alignment
            concerns when optimising proxy rewards; potential misuse of
            fine-tuned models; compute-access disparities that could
            limit replication to well-resourced groups.
      \item \textbf{Environmental.} Estimated 0.4--1.2\,kg
            CO$_2$-equivalent for the Modal H100 experiments; the Tinker
            API footprint is unestimable because the hardware mix is
            undisclosed.
    \end{itemize}}
    \end{enumerate}

% -- 11. Safeguards ------------------------------------------------------------
\item {\bf Safeguards}
    \begin{enumerate}
    \item Does the paper describe safeguards that have been put in place
          for responsible release of data or models that have been
          identified as requiring safeguards?
    \answerNA{The released artefacts are LoRA adapters fine-tuned from
    already-public base models (Qwen, Llama, Nemotron, GPT-OSS,
    Kimi-K2) on standard research datasets (GSM8K, HumanEval, synthetic
    tool-use).  No new high-risk capabilities are introduced relative to
    the base models; released checkpoints inherit the base models'
    licence terms and include standard usage disclaimers in the
    HuggingFace model cards
    (\texttt{huggingface/MODEL\_CARD\_TEMPLATE.md}).}
    \end{enumerate}

% -- 12. Licenses --------------------------------------------------------------
\item {\bf Licenses for existing assets}
    \begin{enumerate}
    \item Are the creators or original owners of assets (e.g., code,
          data, models) used in the paper properly credited, and are
          the licence and terms of use explicitly mentioned and properly
          respected?
    \answerYes{All third-party assets are credited with their licences:
    \begin{itemize}
      \item \textbf{GSM8K}~\cite{cobbe2021training}: MIT licence.
      \item \textbf{HumanEval} (OpenAI): MIT licence.
      \item \textbf{NoRobots} (HuggingFaceH4): CC-BY-NC-4.0 (used for
            the chat-SFT task only).
      \item \textbf{Synthetic tool-use data}: self-generated from public
            API documentation; no upstream licence restrictions.
      \item \textbf{Qwen model family}: Apache-2.0.
      \item \textbf{Llama model family}: Meta Llama Community Licence.
      \item \textbf{Nemotron, GPT-OSS, Kimi-K2}: used under their
            respective published licences; credits in
            Section~\ref{sec:setup}.
      \item \textbf{TRL, PEFT, Transformers}: Apache-2.0 (HuggingFace).
      \item \textbf{Modal}: commercial platform; terms of service
            respected.
      \item \textbf{Tinker API}: proprietary; used under standard API
            terms of service.
      \item \textbf{Our code}: \textbf{Apache-2.0} (see \texttt{LICENSE}
            in the repository root).
    \end{itemize}}
    \end{enumerate}

% -- 13. New Assets ------------------------------------------------------------
\item {\bf New Assets}
    \begin{enumerate}
    \item Are new assets introduced in the paper well documented and is
          the documentation provided alongside the assets?
    \answerYes{New assets and their documentation:
    \begin{itemize}
      \item The \ours{} benchmark suite (harness, evaluation scripts,
            analysis notebooks) at
            \url{https://github.com/arvindcr4/tinker-rl-lab}, released
            under Apache-2.0 and documented by \texttt{README.md},
            \texttt{REPRODUCE.md}, \texttt{ARTIFACT.md},
            \texttt{COMPUTE.md}, \texttt{BASELINES.md},
            \texttt{BENCHMARKS\_COMPARISON.md}, and
            \texttt{LIMITATIONS\_AND\_IMPACT.md}.
      \item LoRA adapter checkpoints from the successful Modal
            experiments at
            \url{https://huggingface.co/arvindcr4/tinker-rl-bench-*},
            each accompanied by a HuggingFace model card generated from
            \texttt{huggingface/MODEL\_CARD\_TEMPLATE.md}
            (intended use, training data, evaluation results, known
            limitations).
      \item All experiment logs on the public Weights~\&~Biases project
            \texttt{arvindcr4-pes-university/tinker-rl-lab-world-class}.
    \end{itemize}
    Checkpoints from JWT-interrupted Tinker runs are \emph{not} uploaded,
    because those jobs did not reach a valid final state
    (Section~\ref{sec:infra_failures}).}
    \end{enumerate}

% -- 14. Crowdsourcing ---------------------------------------------------------
\item {\bf Crowdsourcing and Research with Human Subjects}
    \begin{enumerate}
    \item For crowdsourcing experiments and research with human
          subjects, does the paper include the full text of instructions
          given to participants and screenshots, if applicable, as well
          as details about compensation?
    \answerNA{No crowdsourcing or human subjects research was
    conducted.}
    \end{enumerate}

% -- 15. IRB -------------------------------------------------------------------
\item {\bf Institutional Review Board (IRB) Approvals or Equivalent for
       Research with Human Subjects}
    \begin{enumerate}
    \item Does the paper describe potential risks incurred by study
          participants, whether compensation was provided to study
          participants, and whether informed consent was obtained?
    \answerNA{No human subjects research was conducted; IRB approval is
    not applicable.}
    \end{enumerate}

% -- 16. LLM usage -------------------------------------------------------------
\item {\bf Declaration of LLM usage}
    \begin{enumerate}
    \item Does the paper describe the usage of LLMs if it is an
          important, original, or non-standard component of the core
          methods?
    \answerYes{LLMs appear in this work in two distinct roles, both
    declared:
    \begin{itemize}
      \item \textbf{As subjects of evaluation.} The Qwen, Llama,
            Nemotron, GPT-OSS, and Kimi-K2 model families are the
            primary \emph{subjects} of the benchmark; their role as RL
            fine-tuning targets is described in
            Sections~\ref{sec:benchmark} and~\ref{sec:setup}.
      \item \textbf{As authoring/editing aids.} GitHub Copilot was used
            for boilerplate experiment scripts and LaTeX/BibTeX
            scaffolding.  ChatGPT Pro was used during revision to
            obtain reviewer-style critique of drafts; resulting
            suggestions were evaluated and, where accepted, manually
            incorporated.  All scientific claims, experimental design,
            numerical results, figures, and statistical analyses are
            human-authored and independently verified against the logged
            Weights~\&~Biases traces.
    \end{itemize}
    LLMs are not an \emph{original} or \emph{non-standard} component of
    the core methodology.}
    \end{enumerate}

\end{enumerate}
